\documentclass[11pt,a4paper]{article}
\usepackage[margin=2.2cm]{geometry}
\usepackage[T1]{fontenc}
\usepackage{lmodern}
\usepackage{booktabs}
\usepackage{xcolor}
\usepackage{mdframed}
\usepackage{enumitem}
\usepackage{hyperref}
\hypersetup{colorlinks=true, linkcolor=blue!50!black, urlcolor=blue!50!black}

\definecolor{qbg}{RGB}{232,240,254}
\newmdenv[backgroundcolor=qbg,linewidth=0pt,skipabove=10pt,skipbelow=6pt,
          innertopmargin=6pt,innerbottommargin=6pt]{question}

\newcommand{\Q}[1]{\begin{question}\textbf{#1}\end{question}}

\title{CS683 Checkpoint I --- Viva Preparation\\
\large Microarchitectural Characterisation of Language Runtimes}
\author{Krishna Kukreja \quad 24b3971}
\date{}

\begin{document}
\maketitle
\tableofcontents
\bigskip

\section{Key numbers to have ready}

\begin{center}
\begin{tabular}{lrrr}
\toprule
& \textbf{C} & \textbf{Java} & \textbf{Python} \\
\midrule
IPC, matmul            & 1.738   & 2.086  & 1.461 \\
IPC, list              & 0.04304 & 0.1145 & 1.153 \\
Cycles / multiply-add  & 4.07    & 4.79   & 1{,}513 \\
Cycles / node          & 186     & $\approx$63$^{*}$ & 1{,}279 \\
L2 load MPKI, matmul   & 143.1   & 6.10   & 0.29 \\
LLC load MPKI, list    & 100.6   & 137.3  & 1.83 \\
Branch MPKI, matmul    & 0.279   & 0.391  & 9.11 \\
L1I misses (all windows) & 0     & 0      & 0 \\
\bottomrule
\end{tabular}
\end{center}

\noindent$^{*}$ Distorted --- see Section~\ref{sec:javalist}.

\medskip
\noindent\textbf{Three headline results.}
\begin{itemize}[leftmargin=*,nosep]
  \item Python's time penalty: $372\times$ on matmul, $6.9\times$ on the list walk. Not a constant.
  \item A 16-byte Java array header changes L2 misses by $23\times$ against C.
  \item \texttt{ip\_stride} gives $\le 0.35\%$ anywhere; \texttt{spp\_dev} helps only Python matmul ($+6.7\%$) and hurts Java's list ($-2.8\%$).
\end{itemize}

\section{The tracer patch}

\Q{Walk me through your tracer patch.}

The stock ChampSim Pin tracer has one global \texttt{curr\_instr} buffer, one output
file and shared counters. Pin invokes the analysis routines on every thread, and the
JVM runs 14 of them, so records from compiler and GC threads interleaved into the
same buffer. The result was 23 million records with a zero instruction pointer and
6.8 million with register ids $\ge 128$, which crashed ChampSim with
\texttt{vector::\_M\_range\_check}.

The patch adds a \texttt{-app\_thread N} knob and an early return in four callbacks:
\texttt{ResetCurrentInstruction}, \texttt{BranchOrNot}, \texttt{WriteToSet} and
\texttt{ShouldWrite}. The check is
\texttt{if (PIN\_ThreadId() != KnobAppThread.Value()) return;}

\Q{Why does \texttt{WriteCurrentInstruction} need no check?}

It only executes when \texttt{ShouldWrite} returned true, and \texttt{ShouldWrite}
already filters by thread.

\Q{Why must the check in \texttt{ShouldWrite} come before anything else?}

Because \texttt{ShouldWrite} increments \texttt{skippedInstrCount}. If other threads
reached that line, their instructions would count toward the \texttt{-s} skip and the
trace window would land in the wrong phase of the program.

\Q{Why thread 1 for Java and thread 0 for C and Python?}

Pin numbers threads from 0. The JVM launcher occupies thread 0, so the application
thread is 1. A single-threaded C or Python program is thread 0. The knob defaults to
1, which is why every C and Python trace passes \texttt{-app\_thread 0} explicitly.

\section{Trace windows}

\Q{How did you choose where each window starts?}

C uses ROI marker functions around the kernel. Java and Python use an instruction
skip count measured by a probe run \emph{under the tracer itself}: the patched tracer
prints the application thread's instruction count at exit, and the window is placed
relative to that.

\Q{Why not measure the skip with a lighter tool?}

Because JIT timing depends on tool overhead. The Java application thread executes
9.09\,B instructions under the tracer but 14.58\,B under a light instruction counter.
A skip chosen from the light counter lands in a completely different phase.

\Q{How do you know the windows are in steady state?}

Every workload was traced at two independent positions. All six pairs agree within
3.0\% on IPC, and the sign of every prefetcher effect is identical in both windows.
The largest spread is Python matmul at 3\%, which comes from row-to-row variation in
branch behaviour (MPKI 9.1 vs 7.6) and changes no conclusion.

\Q{What went wrong with one of your early windows?}

A Python list window landed inside interpreter teardown. Branch MPKI was 18.15 and it
reported a \texttt{spp\_dev} speed-up of $+73.8\%$. Moving the window two passes into
the loop, with \texttt{os.\_exit(0)} to skip teardown, gave MPKI 1.50 and $+0.23\%$.

\section{Workload equivalence}

\Q{How do you know the three languages do the same work?}

The matmul checksum is 16{,}392{,}906{,}262{,}816 in all three, and the list checksum
($s \leftarrow 31s + \text{value}$) matches for every pass count used. NumPy was used
only as an external check, never inside a workload.

\Q{Why do your matrix values start at 300?}

CPython caches small integers from $-5$ to 256 as shared objects. Values in that range
would make matrix elements alias to the same object, changing the memory behaviour. An
assertion verifies that all $N^2$ elements have distinct \texttt{id()}.

\Q{Why $N = 512$ and not 256?}

At 256 the whole working set fits comfortably and the L2 effects do not appear. At 512,
B is 1\,MB: it overflows the 512\,KB L2 but fits the 2\,MB LLC, which is the regime that
makes the prefetcher question interesting.

\section{Cache behaviour: the 23\texorpdfstring{$\times$}{x} gap}

\Q{Explain why Java matmul has 6.1 L2 MPKI where C has 143.}

L2 is 1{,}024 sets, 8 ways, 64-byte lines.

\textbf{C:} rows are 2{,}048\,B apart, exactly 32 lines. Each row advances the set index
by 32, and $\gcd(32, 1024) = 32$, so all 512 rows of the column walk of B map into only
32 distinct sets. Those hold $32 \times 8 = 256$ lines, but the walk needs 512. It thrashes.

\textbf{Java:} the 16-byte array header makes rows 2{,}064\,B apart, i.e. 32.25 lines.
Four rows advance 129 lines, and $\gcd(129, 1024) = 1$, so the index cycles through all
1{,}024 sets. The 512-line working set spreads out and fits.

\Q{How do you know the JIT isn't doing a different loop?}

I extracted the source-memory addresses from both traces, grouped them by instruction
pointer, and compared strides.

\begin{center}
\begin{tabular}{lrr}
\toprule
& \textbf{C} & \textbf{Java} \\
\midrule
A load stride & 4\,B & 16\,B \\
B load stride & 2{,}048\,B & 8{,}256\,B \\
Load instructions in loop & 4 & 24 \\
\bottomrule
\end{tabular}
\end{center}

$8{,}256 = 4 \times 2{,}064$ and $16 = 4 \times 4$: C2 unrolled the inner loop four
times, but the traversal order is identical. \texttt{-XX:+PrintCompilation} confirms
\texttt{MatrixMult::multiply} reached tier 4 (C2).

\Q{Does this undermine your "identical work" premise?}

No --- it refines it. The arithmetic is identical, verified by checksum. What differs is
the data layout the runtime imposes, and that is exactly the kind of runtime effect the
project set out to measure.

\section{Prefetchers}

\Q{Why does \texttt{ip\_stride} issue zero prefetches on C matmul?}

It tracks one stride per instruction pointer in a 256-set, 4-way table, needs to observe
the same stride twice before it issues, prefetches a few lines ahead, and stops at a
4\,KB page boundary. B's rows are 2{,}048\,B apart, so only two rows fit in a page and the
stride never repeats inside one. It never confirms, so it never issues. On the list walks
there is no stride at all.

\Q{Why does \texttt{spp\_dev} help Python matmul but barely help C?}

Timeliness, not stride destruction.

\begin{center}
\begin{tabular}{lrrrr}
\toprule
\textbf{spp\_dev} & \textbf{Issued} & \textbf{Useful} & \textbf{Merged} & \textbf{Speed-up} \\
\midrule
C matmul      & 3{,}702{,}262 & 353   & 99.1\% & $+0.58\%$ \\
Python matmul & 241{,}252     & 2{,}540 & $\approx 0\%$ & $+6.70\%$ \\
C list        & 3{,}672{,}718 & 0     & 99.9\% & $-0.01\%$ \\
Java list     & 4{,}766{,}185 & 2     & 54\%   & $-2.80\%$ \\
\bottomrule
\end{tabular}
\end{center}

On C matmul, 99.1\% of requests merge into a demand miss already in flight: right address,
too late. Demand LLC misses fall 74\% and the gain is 0.6\%. On Python there is almost no
merging, so prefetches land before the load needs them and 83\% of demand LLC misses become
hits. The interpreter's long instruction path per data access is what supplies the slack.

\Q{Why does \texttt{spp\_dev} make Java's list walk worse?}

4.77\,M requests, 2.11\,M useless fills, 2 useful. DRAM traffic rises 44\% for no coverage,
and the bandwidth pressure costs about 3\% in cycles.

\Q{Why did you change the prefetch accuracy formula?}

The proposal used \texttt{USEFUL/(USEFUL+USELESS)}. ChampSim's \texttt{USELESS} counter only
increments when an unused prefetched line is evicted, so it undercounts severely. I used
useful $\div$ fills, with fills = L2 prefetch misses $-$ merged, which counts every line
actually brought in.

\section{The Java list artefact}
\label{sec:javalist}

\Q{Java's list IPC is three times better than C's. Isn't that suspicious?}

Yes, and it is an artefact.

\begin{center}
\begin{tabular}{lrr}
\toprule
& \textbf{C list} & \textbf{Java list} \\
\midrule
IPC (no prefetcher)        & 0.0430 & 0.1145 \\
Cycles per LLC load miss   & 231    & 63.6 \\
Average LLC miss latency   & 179    & 213 \\
\bottomrule
\end{tabular}
\end{center}

C completes one miss per 231 cycles, about one miss latency: a true serial chase. Java
completes one per 63.6 cycles although each takes 213, so roughly 3.4 misses overlap. A
strict dependent chain cannot do that. ChampSim tracks dependences through register ids
only, so the JIT-unrolled loads appear independent to it.

Use the Java list for relative prefetcher effects within Java, not for absolute speed
against C. Two reruns with \texttt{-Xbatch} and 300 passes agree (0.1145, 0.1144); earlier
runs with different JIT code gave 0.25 and 0.09, which is why the flags are now fixed.

\section{Front end, branches, TLB}

\Q{Zero L1I misses everywhere --- is that believable?}

In steady state, yes: the hot loop is small in all three runtimes, and CPython 3.12's eval
loop has 140 opcodes (\texttt{len(dis.opmap)}), which fits in 32\,KB. This is consistent with
Rajan et al.'s finding that interpreter I-cache locality is good, though they measured
0.16--1.33\% on much smaller caches, not zero.

The caveat is that JVM \emph{start-up} does miss: 17.1\% of fetches, with 240{,}150 distinct
instruction pointers from class loading and the bytecode interpreter. A single steady-state
window hides this.

\Q{Note on the proposal's opcode claim.}

The proposal said CPython 3.12's bytecode set is "much larger than the $\sim$220 opcodes of
2002". The measurement is 140, i.e.\ fewer. The expectation was wrong and the measurement
stands.

\Q{Where does Python's data-side cost actually sit?}

37--47\% of its L1D accesses are stores (reference counting, stack and frame writes), and
essentially all hit L1D --- store miss rate 0.009--0.024\%. So the cost is instruction count
and L1 traffic, not write misses.

\section{JIT transition}

\begin{center}
\begin{tabular}{lrrr}
\toprule
\textbf{Java matmul} & \textbf{Start-up} & \textbf{300\,M in} & \textbf{Steady} \\
\midrule
IPC                  & 0.727  & 1.559   & 2.086 \\
Branch MPKI          & 12.81  & 0.165   & 0.391 \\
L1I miss rate        & 17.1\% & 0.015\% & 0\% \\
Stores (share)       & 42.3\% & 16.7\%  & 0.05\% \\
\texttt{spp\_dev}    & $+8.8\%$ & $+6.9\%$ & $-0.03\%$ \\
\bottomrule
\end{tabular}
\end{center}

\Q{What does this show?}

IPC triples from start-up to steady state and branch MPKI falls from 12.8 to under 0.4. The
300\,M window is a transition, not steady: compiled code (MPKI 0.17, clean L1I) but different
data behaviour. Note also that \texttt{spp\_dev} helps during warm-up and not in steady state,
so prefetcher conclusions do not transfer between phases.

\section{Limitations}

\begin{itemize}[leftmargin=*]
  \item \textbf{Sampling.} Two windows show stability ($\le 3\%$), not a confidence interval.
  \item \textbf{Java list timing} is distorted (Section~\ref{sec:javalist}).
  \item \textbf{Compiler threads not traced.} Rajan's code-installation writes happen on JIT
        compiler threads, which the per-thread filter excludes. Left for Checkpoint II.
  \item \textbf{One machine model.} Only the L2 prefetcher varies.
  \item \textbf{Prior-work wording.} The proposal's "16\,KB L1" and "no study asks\ldots" claims
        were not verified against the papers. Rohou et al.\ (2015) used Nehalem and Haswell,
        which have 32\,KB L1, so the blanket claim is already wrong for one citation.
\end{itemize}

\Q{What would you do differently?}

More than two windows per workload, so the stability claim becomes a confidence interval. Trace
the JIT compiler threads. And force the Java list into a genuinely serial chain --- for example
by carrying a dependence through memory --- so its timing is faithful.

\section{Checkpoint II}

\begin{enumerate}[leftmargin=*]
  \item Split interpreter overhead from algorithm cost: per-bytecode cost, and a data-versus-instruction breakdown.
  \item Establish the mechanism behind \texttt{spp\_dev}'s Python gain and Java-list failure.
  \item Implement one targeted mitigation in ChampSim and report recovered IPC and MPKI.
\end{enumerate}

\vspace{1em}
\noindent\rule{\textwidth}{0.4pt}
\vspace{0.5em}

\noindent\small Repository: \url{https://github.com/Krishna-get/cs683-language-runtimes}

\end{document}